% ============================================================
\chapter{\'Equation de Laplace et Fonctions Harmoniques}
\label{ch:laplace}
% ============================================================

L'\'equation de Laplace $\Delta u = 0$ est la reine des EDP elliptiques. Ses solutions, les \emph{fonctions harmoniques}, mod\'elisent les \'etats d'\'equilibre~: temp\'erature stationnaire, potentiel \'electrostatique, \'ecoulement irrotationnel. Pierre-Simon Laplace l'a \'etudi\'ee d\`es 1782 dans le contexte de la gravitation newtonienne, et ses propri\'et\'es --- principe du maximum, r\'egularit\'e infinie, formule de repr\'esentation de Poisson --- en font l'une des \'equations les mieux comprises de toute la physique math\'ematique. Ce chapitre d\'eveloppe ces propri\'et\'es remarquables.

\section{Introduction et motivation}

L'\'equation de Laplace $\Delta u = 0$ mod\'elise les \'etats d'\'equilibre~:
temp\'erature stationnaire, potentiel \'electrostatique en l'absence de
charges, \'ecoulement potentiel. Son \'etude r\'ev\`ele des propri\'et\'es
remarquables de r\'egularit\'e, de valeur moyenne et de maximum.

\begin{definition}[Fonction harmonique]
Une fonction $u \in C^2(\Omega)$ est \textbf{harmonique} dans un ouvert
$\Omega \subset \R^n$ si~:
\begin{equation}\label{eq:laplace}
    \Delta u = \sum_{i=1}^{n} \frac{\partial^2 u}{\partial x_i^2} = 0
    \quad \text{dans } \Omega.
\end{equation}
\end{definition}

\begin{definition}[\'Equation de Poisson]
L'\textbf{\'equation de Poisson} est l'\'equation non homog\`ene~:
\begin{equation}\label{eq:poisson}
    -\Delta u = f \quad \text{dans } \Omega.
\end{equation}
\end{definition}

\section{Solution fondamentale}

\begin{theoreme}[Solution fondamentale du laplacien]
La \textbf{solution fondamentale} de l'\'equation de Laplace en dimension $n$
est~:
\begin{equation}\label{eq:sol_fondamentale_laplace}
    \Gamma(x) = \begin{cases}
        -\frac{1}{2\pi}\ln\abs{x} & \text{si } n = 2,\\[6pt]
        \frac{1}{n(n-2)\omega_n}\frac{1}{\abs{x}^{n-2}} & \text{si } n \geq 3,
    \end{cases}
\end{equation}
o\`u $\omega_n = \frac{\pi^{n/2}}{\Gamma(n/2+1)}$ est le volume de la boule
unit\'e de $\R^n$. Elle v\'erifie $-\Delta\Gamma = \delta_0$ au sens des
distributions.
\end{theoreme}

\section{Propri\'et\'e de la valeur moyenne}

\begin{theoreme}[Propri\'et\'e de la valeur moyenne]
Si $u$ est harmonique dans $\Omega$ et $B(x,r) \subset \Omega$, alors~:
\begin{equation}\label{eq:valeur_moyenne_sphere}
    u(x) = \frac{1}{n\omega_n r^{n-1}}\int_{\partial B(x,r)} u(y)\, dS(y)
    \quad \text{(moyenne sur la sph\`ere)},
\end{equation}
\begin{equation}\label{eq:valeur_moyenne_boule}
    u(x) = \frac{1}{\omega_n r^n}\int_{B(x,r)} u(y)\, dy
    \quad \text{(moyenne sur la boule)}.
\end{equation}
\end{theoreme}

\begin{proof}
Posons $\phi(r) = \frac{1}{n\omega_n r^{n-1}}\int_{\partial B(x,r)} u\, dS$.
Par le changement $y = x + r\sigma$, $\sigma \in \partial B(0,1)$~:
\[
    \phi(r) = \frac{1}{n\omega_n}\int_{\partial B(0,1)} u(x+r\sigma)\, dS(\sigma).
\]
D\'erivons~:
\[
    \phi'(r) = \frac{1}{n\omega_n}\int_{\partial B(0,1)}
    \nabla u(x+r\sigma) \cdot \sigma\, dS(\sigma)
    = \frac{1}{n\omega_n r^{n-1}}\int_{\partial B(x,r)}
    \frac{\partial u}{\partial \nu}\, dS.
\]
Par le th\'eor\`eme de la divergence et $\Delta u = 0$~:
\[
    \phi'(r) = \frac{1}{n\omega_n r^{n-1}}\int_{B(x,r)} \Delta u\, dy = 0.
\]
Donc $\phi(r) = \phi(0^+) = u(x)$ par continuit\'e. La formule sur la
boule s'obtient en int\'egrant en $r$.
\end{proof}

\begin{intuition}
La valeur d'une fonction harmonique en un point est la moyenne de ses
valeurs sur toute sph\`ere centr\'ee en ce point. C'est l'analogue continu
du fait qu'en \'equilibre thermique, la temp\'erature en un point est la
moyenne des temp\'eratures environnantes.
\end{intuition}

\begin{theoreme}[R\'eciproque de la valeur moyenne]
Si $u \in C(\Omega)$ v\'erifie la propri\'et\'e de la valeur moyenne
\eqref{eq:valeur_moyenne_sphere} pour toute boule $B(x,r) \subset \Omega$,
alors $u$ est harmonique dans $\Omega$ (et en particulier $C^\infty$).
\end{theoreme}

\section{Principe du maximum}

\begin{theoreme}[Principe du maximum fort pour les fonctions harmoniques]
\label{thm:max_fort_harmonique}
Soit $\Omega \subset \R^n$ un ouvert connexe born\'e et $u \in C^2(\Omega) \cap C(\overline{\Omega})$
harmonique dans $\Omega$. Alors~:
\begin{enumerate}
    \item $\max_{\overline{\Omega}} u = \max_{\partial\Omega} u$
          (principe du maximum).
    \item Si le maximum est atteint en un point int\'erieur, alors $u$ est
          constante dans $\Omega$ (principe du maximum fort).
\end{enumerate}
De m\^eme pour le minimum.
\end{theoreme}

\begin{proof}
Supposons que $u$ atteint son maximum $M$ en $x_0 \in \Omega$. Soit
$r > 0$ tel que $B(x_0, r) \subset \Omega$. Par la valeur moyenne~:
\[
    M = u(x_0) = \frac{1}{\omega_n r^n}\int_{B(x_0,r)} u(y)\, dy.
\]
Comme $u(y) \leq M$ partout et $u(x_0) = M$, la moyenne ne peut valoir $M$
que si $u \equiv M$ dans $B(x_0,r)$. L'ensemble $\{x : u(x) = M\}$ est
donc ouvert (et ferm\'e dans $\Omega$ par continuit\'e). Par connexit\'e,
$u \equiv M$ dans $\Omega$.
\end{proof}

\begin{corollaire}[Unicit\'e du probl\`eme de Dirichlet]
Le probl\`eme de Dirichlet $\Delta u = 0$ dans $\Omega$, $u = g$ sur
$\partial\Omega$, admet au plus une solution.
\end{corollaire}

\begin{corollaire}[Stabilit\'e]
Si $u$ et $v$ sont harmoniques avec $u = g_1$, $v = g_2$ sur
$\partial\Omega$, alors~:
\[
    \max_{\overline{\Omega}} \abs{u - v} \leq \max_{\partial\Omega} \abs{g_1 - g_2}.
\]
\end{corollaire}

\section{R\'egularit\'e des fonctions harmoniques}

\begin{theoreme}[R\'egularit\'e $C^\infty$]
\label{thm:regularite_harmonique}
Toute fonction harmonique est $C^\infty$. Plus pr\'ecis\'ement, si $u$
est harmonique dans $\Omega$, alors $u \in C^\infty(\Omega)$ et
toutes ses d\'eriv\'ees partielles sont harmoniques.
\end{theoreme}

\begin{theoreme}[Analyticit\'e]
Toute fonction harmonique est analytique r\'eelle~: elle est \'egale \`a
sa s\'erie de Taylor en tout point de son domaine.
\end{theoreme}

\begin{theoreme}[Estimations de gradient]
Si $u$ est harmonique dans $B(x_0, R)$, alors~:
\begin{equation}\label{eq:estimation_gradient}
    \abs{\nabla u(x_0)} \leq \frac{n}{R}\max_{\partial B(x_0,R)} \abs{u}.
\end{equation}
Plus g\'en\'eralement, pour tout multi-indice $\alpha$~:
\[
    \abs{D^\alpha u(x_0)} \leq \frac{C_{n,\abs{\alpha}}}{R^{\abs{\alpha}}}
    \max_{\partial B(x_0,R)} \abs{u}.
\]
\end{theoreme}

\section{In\'egalit\'e de Harnack}

\begin{theoreme}[In\'egalit\'e de Harnack]
Soit $u$ harmonique et \textbf{positive} dans $B(x_0, R) \subset \R^n$.
Pour tout $x \in B(x_0, R)$ avec $\abs{x - x_0} = r < R$~:
\begin{equation}\label{eq:harnack}
    \frac{R^{n-2}(R-r)}{(R+r)^{n-1}}\, u(x_0) \leq u(x)
    \leq \frac{R^{n-2}(R+r)}{(R-r)^{n-1}}\, u(x_0).
\end{equation}
En particulier, sur tout compact $K \Subset \Omega$, il existe $C = C(K, \Omega)$
tel que~:
\[
    \max_K u \leq C \min_K u.
\]
\end{theoreme}

\begin{remarque}
L'in\'egalit\'e de Harnack affirme que les fonctions harmoniques positives
ne peuvent pas varier trop rapidement. Elle est fondamentale en th\'eorie
de la r\'egularit\'e.
\end{remarque}

\section{Formule de repr\'esentation de Green}

\begin{definition}[Fonction de Green]
La \textbf{fonction de Green} du domaine $\Omega$ est $G(x,y) = \Gamma(x-y) - h^y(x)$,
o\`u $h^y$ est harmonique dans $\Omega$ avec $h^y = \Gamma(\cdot - y)$ sur
$\partial\Omega$.
\end{definition}

\begin{theoreme}[Formule de repr\'esentation]
Si $u \in C^2(\overline{\Omega})$ est harmonique dans $\Omega$, alors
pour tout $x \in \Omega$~:
\begin{equation}\label{eq:representation_green}
    u(x) = -\int_{\partial\Omega} u(y)\, \frac{\partial G}{\partial \nu_y}(x,y)\, dS(y).
\end{equation}
Pour l'\'equation de Poisson $-\Delta u = f$~:
\begin{equation}\label{eq:representation_poisson}
    u(x) = \int_\Omega G(x,y)\, f(y)\, dy
    - \int_{\partial\Omega} u(y)\, \frac{\partial G}{\partial \nu_y}(x,y)\, dS(y).
\end{equation}
\end{theoreme}

\begin{exemple}[Fonction de Green du disque]
Pour le disque $B(0,R) \subset \R^2$, la fonction de Green est~:
\[
    G(x,y) = -\frac{1}{2\pi}\ln\abs{x-y}
    + \frac{1}{2\pi}\ln\!\left(\frac{\abs{x}\abs{y}}{R}\abs{\frac{x}{\abs{x}^2} \cdot R^2 - y}\right).
\]
La formule de repr\'esentation redonne le \textbf{noyau de Poisson}.
\end{exemple}

\section{Th\'eor\`eme de Liouville}

\begin{theoreme}[Th\'eor\`eme de Liouville]
Toute fonction harmonique born\'ee sur $\R^n$ est constante.
\end{theoreme}

\begin{proof}
Soit $u$ harmonique born\'ee sur $\R^n$, $\abs{u} \leq M$. Pour tout
$x_0 \in \R^n$ et $R > 0$, l'estimation du gradient donne~:
\[
    \abs{\nabla u(x_0)} \leq \frac{n}{R}\, M.
\]
En faisant $R \to \infty$, on obtient $\nabla u(x_0) = 0$ pour tout $x_0$,
donc $u$ est constante.
\end{proof}

\section{Exercices}

\begin{exercice}
Montrer que si $u$ est harmonique dans $B(0,R) \setminus \{0\} \subset \R^2$
et born\'ee, alors la singularit\'e en $0$ est \'eliminable (th\'eor\`eme
de singularit\'e \'eliminable).
\end{exercice}

\begin{exercice}
Utiliser la formule de Poisson pour r\'esoudre le probl\`eme de Dirichlet
dans le disque $B(0,1)$ avec $u(1,\theta) = \cos(2\theta) + 3\sin(\theta)$.
\end{exercice}

\begin{exercice}
Montrer que si $u$ est harmonique et positive dans $\R^n$, alors $u$
est constante (cons\'equence de Harnack).
\end{exercice}

\begin{exercice}
Calculer la fonction de Green du demi-espace $\{x_n > 0\}$ par la
m\'ethode des images.
\end{exercice}

\begin{exercice}
Montrer que la valeur moyenne d'une fonction harmonique sur une sph\`ere
est \'egale \`a sa valeur au centre. Que se passe-t-il pour les fonctions
sous-harmoniques ($\Delta u \geq 0$)~?
\end{exercice}

\begin{exercice}
Soit $u$ harmonique dans $\Omega = B(0,1) \setminus \overline{B(0,1/2)}
\subset \R^3$, avec $u = 0$ sur $\abs{x} = 1/2$ et $u = 1$ sur $\abs{x} = 1$.
Trouver $u$ sous la forme $u(r) = a + b/r$.
\end{exercice}

\begin{exercice}[Principe de Schwarz]
Soit $u$ harmonique dans le demi-disque $\{r < 1, 0 < \theta < \pi\}$
avec $u = 0$ sur le diam\`etre $[-1,1]$. Montrer que $u$ se prolonge en
une fonction harmonique dans tout le disque par r\'eflexion impaire.
\end{exercice}
